\section{Ruta de implementación propuesta}

La evolución debe ser incremental para conservar el comportamiento actual. Cada etapa puede integrarse y probarse de manera independiente.

\subsection{Etapa 1: consolidar contratos y nombres}

\begin{enumerate}
  \item Definir una convención: tipos en \code{PascalCase}, campos serializados en \code{camelCase} y acciones con verbo explícito.
  \item Resolver términos ambiguos del diagrama: \code{skinId}, \code{imageId}, \code{themeId}, \code{dimensions}, \code{colliderSize} y \code{colliderOffset}.
  \item Definir IDs como \code{string} si pueden provenir de servicios o bases de datos heterogéneas.
  \item Crear contratos de inventario, interacción y área antes de agregarlos a los estados principales.
\end{enumerate}

\subsection{Etapa 2: completar PlayerState}

Agregar identidad y apariencia al modelo existente sin reemplazar la ruta de movimiento. Cada dato nuevo requiere, como mínimo:

\begin{itemize}
  \item campo privado serializado y propiedad de solo lectura;
  \item parámetro de constructor con valor seguro;
  \item método \code{With...};
  \item tipo de acción, payload y fábrica de acción;
  \item caso en el reductor;
  \item comparación en \code{StatesAreEqual};
  \item prueba del reductor y, si produce un efecto visual, prueba del adaptador.
\end{itemize}

El correo electrónico debe permanecer en el estado de cuenta o perfil si no es necesario para la simulación del jugador. Separarlo reduce exposición accidental y evita replicar datos privados a otros clientes.

\subsection{Etapa 3: implementar WorldState}

Crear el módulo \code{World/State} con \code{WorldStateData}, \code{WorldAction}, \code{WorldStateReducer} y \code{WorldStateStore}. Su primera versión puede manejar mundo, sala, tema y dimensiones. La carga completa debe ser atómica:

\begin{lstlisting}[caption={Ejemplo de acción compuesta para cargar un mundo}]
public static WorldAction LoadWorld(WorldDefinition definition) =>
    new WorldAction(
        WorldActionType.LoadWorld,
        new WorldPayload(
            definition.Id,
            definition.DisplayName,
            definition.RoomId,
            definition.ThemeId,
            definition.Dimensions));
\end{lstlisting}

\subsection{Etapa 4: implementar decoraciones como colección}

Un único \code{ObjDecorationStateData} describe una instancia. El store de mundo debería custodiar una colección o un diccionario por ID y ofrecer acciones como:

\begin{itemize}
  \item \code{AddDecoration(payload)};
  \item \code{RemoveDecoration(id)};
  \item \code{MoveDecoration(id, position)};
  \item \code{SetDecorationTheme(id, themeId)};
  \item \code{SetDecorationInventory(id, inventory)}.
\end{itemize}

Para actualizaciones frecuentes, un diccionario facilita buscar por ID. Para serialización con el \code{JsonUtility} tradicional de Unity, una lista es más directa; puede mantenerse una lista serializable y construir un índice de runtime no serializado.

\subsection{Etapa 5: fachada de estado principal}

Solo cuando un caso real lo requiera, crear una fachada que exponga una instantánea coherente de sesión, mundo, jugador y decoraciones. No es obligatorio fusionar todos los stores en un único \code{MonoBehaviour}. Una fachada de solo lectura puede componerlos:

\begin{lstlisting}[caption={Fachada conceptual de solo lectura}]
public readonly struct PrincipalStateSnapshot
{
    public GameSessionStateData Session { get; }
    public WorldStateData World { get; }
    public PlayerStateData Player { get; }
    public IReadOnlyList<ObjDecorationStateData> Decorations { get; }
}
\end{lstlisting}

Esta opción conserva la independencia de cada ciclo de actualización y permite producir guardados, telemetría o depuración global sin habilitar mutaciones externas.

\subsection{Mapa completo de referencia}

La Figura~\ref{fig:mapa-completo} reúne los estados y los métodos previstos. Debido a su relación de aspecto, se presenta en una página horizontal para conservar legibilidad.

\begin{landscape}
\begin{figure}[p]
  \centering
  \includegraphics[width=0.98\linewidth,height=0.86\textheight,keepaspectratio]{\figurepath/mapa-estados-metodos.png}
  \caption{Mapa consolidado de estados, propiedades y operaciones planeadas.}
  \label{fig:mapa-completo}
\end{figure}
\end{landscape}

Los métodos del mapa deben interpretarse como inventario de capacidades. En la implementación, los getters son propiedades de solo lectura y los setters se convierten en acciones; no se recomienda copiar literalmente pares \code{set/get} para cada campo.
